\chapter{Discussion and
Conclusion}\label{chapter-6-discussion-and-conclusion}

\section{Overview}\label{60-overview}

This chapter synthesizes the empirical findings of Chapter 5 against the
four plan contributions, situates the work within the
cross-simulator benchmarking literature, catalogs the limitations that
bound the claims, identifies concrete future-work directions, and closes
with the central thesis statement that SimForge supports.

Chapter 3 (Methods) presented \emph{how} the framework is built. Chapter
4 (Experiments) described \emph{what} was measured. Chapter 5 (Results)
reported the measurements. This chapter answers \emph{so what}: which of
the proposed research questions were answered, which only partially,
where the framework\textquotesingle s adapter pattern proved more
general than expected (the LPSim → DTALite substitution, the
cross-platform reproducibility analysis), and which threads remain open
for the research community to pick up.

The headline takeaway: the plan-§1.10 Objective of \emph{"a
reproducible, transparent, and extensible cross-simulator testing
framework"} is delivered, with four empirical findings that emerged from
the build and now form part of the framework\textquotesingle s
contribution beyond the original C1-C4 list.

\begin{center}\rule{0.5\linewidth}{0.5pt}\end{center}

\section{Contributions revisited}\label{61-contributions-revisited}

The December 2025 plan committed to four principal contributions
(§1.11 C1-C4). Each is reviewed below with what was actually delivered,
along with the headline measurement that validates it.

\subsection{C1 --- Canonical schema and
validators}\label{611-c1--canonical-schema-and-validators}

\textbf{Plan commitment}: \emph{"A mixed-format, simulator-agnostic
input bundle consisting of \texttt{network.xml}, \texttt{demand.csv},
\texttt{signals.xml}, \texttt{config.xml}, and a manifest file. The
schema enforces field-level validation rules and pre-run integrity
checks."} (§1.11 C1)

\textbf{Delivered}: the canonical 5-file bundle structure ships in
\texttt{scenarios/\textless{}scenario\_id\textgreater{}/} for every
scenario. Field-level validation is implemented in
\texttt{pipeline/validation/validate\_bundle.py}, exercised by every
adapter at run time, and covered by \textasciitilde626 tests across 31
test files (Chapter 3 §3.2.8, Chapter 4 §4.7). The
bundle\textquotesingle s \texttt{manifest.xml} records SHA-256 hashes
for downstream provenance verification.

The cross-engine fairness contract (Q1 byte-identity, §5.6) \textbf{is
the empirical proof} that C1 is delivered. If every
engine\textquotesingle s \texttt{feasibility\_report.json} is
byte-identical across the three engines, the canonical schema is doing
its job --- converting one set of inputs into bit-stable engine-specific
inputs.

\subsection{C2 --- Deterministic
adapters}\label{612-c2--deterministic-adapters}

\textbf{Plan commitment}: \emph{"One-to-one, auditable mappings from
the canonical schema to each simulator\textquotesingle s native input
structure for SUMO, MATSim, POLARIS, LPSim, and QarSUMO, ensuring
consistent parameter translation, unit normalization, and time-step
synchronization."} (§1.11 C2)

\textbf{Delivered (with documented engine substitution)}: three engines
ship --- SUMO, MATSim, and DTALite. The original LPSim and QarSUMO did
not survive integration; POLARIS was license-gated and ruled out at the
criteria stage. The substitution rationale is documented in
\texttt{doc/engines/LPSIM\_RETROSPECTIVE.md},
\texttt{doc/engines/QARSUMO\_RETROSPECTIVE.md}, and
\texttt{doc/engines/THIRD\_ENGINE\_OPTIONS.md}. The empirical
implication is preserved: the three shipped engines cover three distinct
paradigms --- microscopic queue (SUMO meso), activity-based queue
(MATSim qsim), and user-equilibrium DTA (DTALite) --- which is
sufficient for the cross-paradigm comparison the plan needed.

Within the shipped roster, all three adapters are byte-deterministic
under fixed seeds (verified by the \texttt{determinism} test marker),
all three share the SCC + feasibility filter
(\texttt{adapters/common/feasibility.py}), and all three are
mutation-tested (\texttt{doc/MUTATION\_BASELINE.md}).

The substitution is itself a contribution of the adapter pattern: the
framework absorbed three engine ruleouts (LPSim, QarSUMO, POLARIS) and
one swap-in (DTALite) without disturbing C1, C3, or C4. The contribution
is therefore \emph{engine-paradigm coverage}, not
\emph{every-named-engine-integrated}.

\subsection{C3 --- Pinned-digest reproducible
execution}\label{613-c3--pinned-digest-reproducible-execution}

\textbf{Plan commitment}: \emph{"Implementation of version-pinned,
containerized workflows (OCI/Singularity) with fixed random seeds,
consistent time-step policies, and hash-verified inputs and outputs."}
(§1.11 C3)

\textbf{Delivered (Wave 2, 2026-05-20)}: a \texttt{Dockerfile} at repo
root, a GitHub Actions workflow
(\texttt{.github/workflows/build-container.yml}) that auto-builds and
publishes to GHCR
(\texttt{ghcr.io/phanidharakula/simforge:\textless{}git-sha\textgreater{}}),
a pinned-digest manifest at \texttt{lib/container/manifest.json}, and
opt-in container mode in the canonical SBATCH wrappers via the
\texttt{SIMFORGE\_USE\_CONTAINER=1} environment variable. The container
was end-to-end verified on OSC Cardinal with Apptainer 1.4.5 on
2026-05-20.

The Wave 2 work surfaced a finding that elevates C3 from "deliverable"
to "novel methodological contribution": empirically measured
\emph{cross-platform reproducibility limits} of approximately 0.2-3 \%
per engine, documented in §5.6.3. This is discussed separately in §6.2.3
below.

\subsection{C4 --- Hardware-normalized
evaluation}\label{614-c4--hardware-normalized-evaluation}

\textbf{Plan commitment}: \emph{"Definition of key performance
indicators (KPIs) for fidelity (RMSE, GEH, KS statistics), scalability
(runtime vs agents, vehicles/sec per core, vehicles/sec per watt), and
reproducibility (R = 1 − σ/μ), all reported with confidence intervals
and hardware normalization."} (§1.11 C4)

\textbf{Delivered (with two documented adjustments)}: the
reproducibility index R = 1 − σ/μ is implemented at
\texttt{evaluation/metrics/reproducibility.py} and reported on every
cell (Table 5.2). 95 \% confidence-interval half-widths on every mean
are implemented at \texttt{evaluation/metrics/confidence.py}
(Student\textquotesingle s t) and appear in every results table. The
cross-engine Q1-Q5 fairness audit is at
\texttt{evaluation/audit\_fairness.py}. The reproducibility scorecard is
at \texttt{tools/generate\_scorecard.py} (auto-emitted on every
\texttt{run\_benchmark} invocation).

The two documented adjustments from the plan text:

\begin{enumerate}
\def\labelenumi{\arabic{enumi}.}
\item
  \textbf{Fidelity metrics RMSE/GEH/KS are implemented but not wired to
  observed baselines.} The functions exist at
  \texttt{evaluation/metrics/fidelity.py} but the ingestion pipeline for
  Chicago Traffic Counts / LADOT / NYC DOT loop-detector data is not
  built. The shipped substitute for fidelity-vs-reality is the
  cross-engine Q4 mean-TT comparison, framed as paradigm divergence
  rather than fidelity error.
\item
  \textbf{Per-watt normalization is not reported.} OSC HPC nodes do not
  expose live power meters at the job-allocation level (RAPL counters
  are root-only). The wall-time + per-core normalization is shipped;
  per-watt is omitted.
\end{enumerate}

Both adjustments are catalogued in \texttt{doc/PLAN\_DEVIATIONS.md}
(D2, I1) and discussed in §6.4 below.

\subsection{Summary table}\label{615-summary-table}

\begin{longtable}[]{@{}llll@{}}
\toprule\noalign{}
Contribution & Status & Empirical evidence & Adjustments \\
\midrule\noalign{}
\endhead
\bottomrule\noalign{}
\endlastfoot
C1 schema + validators & ✓ Delivered & Q1 PASS at every tier (5
scenarios × 3 engines) & None \\
C2 deterministic adapters & ✓ Delivered (3 of 5 engines) & Determinism
test marker; R = 1.0 within context for MATSim/DTALite &
LPSim/QarSUMO/POLARIS ruled out, DTALite substituted \\
C3 pinned-digest container & ✓ Delivered (Wave 2 hotfix, 2026-05-20) &
Container verified on Cardinal; cross-platform measurement & Wave 2 was
post-original-timeline \\
C4 KPIs with CIs & ✓ Delivered (R, 95 \% CI, audit, scorecard) & Tables
5.1 + 5.2, audit\_fairness.txt, scorecard.md & Fidelity not wired to
observed baselines; per-watt dropped \\
\end{longtable}

\begin{center}\rule{0.5\linewidth}{0.5pt}\end{center}

\section{Synthesis of empirical
findings}\label{62-synthesis-of-empirical-findings}

Beyond the four C-rows, three headline empirical findings emerged during
the build (§6.2.1-6.2.3) that contribute to the cross-simulator
benchmarking literature in their own right, plus a within-engine sibling
result (§6.2.4) that reinforces the cross-engine paradigm finding. They
are not new C-rows, but they are not incidental either --- each surfaces
a measurement the cross-simulator literature does not typically report.

\pandocbounded{\includegraphics[keepaspectratio]{../figures/fig_6_1_paradigm_spread.png}}

\subsection{Phase 14: BFS deduplication makes the large tier
tractable}\label{621-phase-14-bfs-deduplication-makes-the-large-tier-tractable}

\textbf{Finding}: pre-Phase-14, the chicago\_200k\_car benchmark on
Cardinal required 141.87 h wall (job 9332478). Post-Phase-14 (cold
cache): \textasciitilde7.14 h wall. Warm-cache re-runs:
\textasciitilde37 min. The compounding speedups are \textasciitilde20 ×
cold-vs-cold and \textasciitilde228 × for warm-cache re-runs against the
Phase 13 baseline.

\textbf{Mechanism} (Chapter 3 §3.8): the SUMO and MATSim adapters each
ran their own per-trip BFS pass through the SCC, paying the full routing
cost twice per scenario. Phase 14 introduces a shared
\texttt{canonical\_routes} module
(\texttt{adapters/common/canonical\_routes.py}) that computes the BFS
pass once per scenario, parallel-multiprocessing across all available
SLURM cores, and feeds the resulting routes to both adapters. Phase
14.12 further collapsed MATSim\textquotesingle s per-trip
\texttt{find\_link\_for\_*} from O(N) linear scan to O(1) indexed
lookup. Phase 14.13 hoisted the cache file from a per-output-dir
location to a global, content-addressable
\texttt{cache/canonical\_routes/canonical\_routes\_\textless{}sha256\textgreater{}.jsonl}
so that all runs of the same scenario share one cache regardless of
which \texttt{-\/-output} directory invokes the BFS.

\textbf{Implication}: the plan\textquotesingle s largest declared
tier (5 M trips, §1.13) remained out of compute budget, but the
second-largest (500 K, nyc\_500k\_car) became tractable for the first
time --- completed in 12 h 8 min wall (job 9971042, 2026-05-19). The 200
K tier moved from a 6-day cold-cache run to a 37-min warm re-run, which
is the difference between a \emph{demonstration} result and an
\emph{iterable} result. The framework can now support systematic
parameter sweeps at the 200 K-trip metropolitan tier.

\subsection{Q4 paradigm divergence at
scale}\label{622-q4-paradigm-divergence-at-scale}

\textbf{Finding} (Chapter 5 §5.6.2): the SUMO meso vs MATSim qsim
mean-TT ratio is not regime-stable. On the small tier (1 K - 50 K trips)
the ratio narrows from 0.869 at chicago\_1k\_car to 1.046 at
la\_50k\_car, within ±5 \%. On the large tier this convergence sharply
reverses: chicago\_200k\_car ratio is 0.645 (−35.5 \%), and
nyc\_500k\_car ratio is 0.037 (−96.3 \%). The latter is the sharpest
paradigm divergence in the dataset and is documented as the headline
finding of §5.6.2.

\textbf{Mechanism}: SUMO meso refuses vehicle insertion at congested
origin-edges; vehicles that cannot insert are dropped from the
simulation and never appear in \texttt{tripinfo.xml}. MATSim qsim holds
vehicles in queue until they can advance; the wait time is counted
toward the trip\textquotesingle s travel time. At low congestion both
engines report on nearly the same trip set and report similar means; at
high congestion SUMO\textquotesingle s reported mean is biased toward
the "easier" subset that did insert, while MATSim\textquotesingle s mean
accumulates multi-hour queue waits. NYC\textquotesingle s
bridge-and-corridor topology concentrates flow more than
Chicago\textquotesingle s grid, so the saturation is sharper at the same
per-trip density.

\textbf{Significance}: the literature commonly assumes that cross-engine
mean-TT comparisons become more reliable at scale via the
law-of-large-numbers smoothing of per-trip variance. The SimForge
measurement shows the opposite: at scale, the \emph{paradigm} signal
becomes more pronounced, not less, because saturation engages the
engines\textquotesingle{} divergent congestion-handling behaviour. This
is a methodological caution for any cross-simulator study reporting
single-engine TTs at high congestion density.

The fairness contract (Q1-Q3 PASS at chicago\_200k\_car and
nyc\_500k\_car) is what makes this finding \emph{interpretable}. Without
byte-identical feasibility verdicts, byte-identical SCC networks, and
identical simulated trip-count targets, the −96.3 \% gap could be
attributed to input asymmetry. With Q1-Q3 PASS empirically demonstrated,
the gap is \emph{necessarily} paradigm-attributable.

\subsection{Cross-platform reproducibility
limits}\label{623-cross-platform-reproducibility-limits}

\pandocbounded{\includegraphics[keepaspectratio]{../figures/fig_6_2_reproducibility_regimes.png}}

\textbf{Finding} (Chapter 5 §5.6.3): at a fixed code version, SimForge
is bit-identical across every measured platform. The 0.2-3 \% shifts
originally reported in §5.6.3 (2.95 \% MATSim, 0.24 \% DTALite) were
\textbf{code-version drift between time-separated comparisons} (Pitzer
Phase 12 ↔ Cardinal Phase 14.13), not cross-platform variability. Two
follow-up measurements isolated the axes:

\begin{itemize}
\tightlist
\item
  \textbf{Cross-execution-context on x86\_64} (§5.6.3.1): Cardinal RHEL
  host venv with Adoptium OpenJDK 21 ↔ Cardinal Debian container with
  apt OpenJDK 17. 20/20 cells at chicago\_200k\_car + nyc\_500k\_car
  produce byte-identical SUMO + MATSim outputs.
\item
  \textbf{Cross-architecture} (§5.6.3.2): Mac arm64 (Apple OpenJDK 17) ↔
  Cardinal Sapphire Rapids x86\_64. MATSim mean TT 318.774 ↔ 318.77 (4
  sig figs); DTALite mean TT 172.5605 ↔ 172.56 (4 sig figs). Consistent
  with bit-identity at the precision of the cited Cardinal numbers.
\end{itemize}

The original §5.6.3 attribution to "JVM build differences" (refuted by
§5.6.3.1) and then to "CPU architecture" (refuted by §5.6.3.2) was wrong
on both counts. The actual cause is the Phase 14 canonical\_routes BFS
replacing the legacy per-adapter BFS in
\texttt{build\_matsim\_plans\_xml}, producing slightly different
per-trip route paths that propagate to mean-TT drift.

\textbf{Significance}: cross-platform variability at fixed code is not
the reproducibility risk the cross-simulator benchmarking literature
typically assumes. Code-version drift on the same platform is. The
pinned-digest container is load-bearing because it freezes the code +
bundle + toolchain at a precise git SHA --- not because it equalizes
platform differences (which are already minimal).

\textbf{Actionable advice for the cross-simulator benchmarking
community: cite the container digest, not the version number.} A version
pin like \emph{"version 1.26.0"} allows code-version drift if any
dependent code (routing layer, bundle generator, adapter glue) changes
between measurements; the digest pins everything at a precise moment.
The fairness contract (Q1 byte-identity, Q2 same SCC, Q3 same trip
count) is platform-invariant because it operates on pre-engine inputs
that run in pure Python on identical canonical data.

\subsection{Within-engine paradigm spread: SUMO micro vs meso
converges at
saturation}\label{624-within-engine-paradigm-spread-sumo-micro-vs-meso-converges-at-saturation}

\textbf{Finding} (Chapter 5 §5.4, extended with the chicago\_200k\_car
SUMO micro pilot, Cardinal job 9980007, 2026-05-19/20): within a single
engine, the micro vs meso mean-TT gap \emph{narrows} at saturation
density, while the wall-time premium \emph{grows monotonically}. The
pilot extended the within-engine SUMO comparison from the small tier (1K
+ 10K) to the large tier (200K), and the trajectory inverted:

\begin{longtable}[]{@{}lrrrr@{}}
\toprule\noalign{}
Tier & meso TT & micro TT & mean-TT Δ & wall ratio \\
\midrule\noalign{}
\endhead
\bottomrule\noalign{}
\endlastfoot
chicago\_1k\_car & 268.93 s & 342.14 s & +27.2 \% & 1.65 × \\
nyc\_10k\_car & 661.6 s & 1143.8 s & +72.9 \% & 9.01 × \\
\textbf{chicago\_200k\_car} & \textbf{8746.40 s} & \textbf{9430.04 s} &
\textbf{+7.8 \%} & \textbf{\textasciitilde124 ×} \\
\end{longtable}

At 1K and 10K the micro/meso TT gap widens with scale, consistent with
the intuition that microscopic dynamics capture intersection delays and
queue spillback that mesoscopic queue-based mobsim averages out --- and
that this captured detail accumulates as congestion grows. The 200K data
point falsifies that intuition: at saturation, the gap narrows to +7.8
\%.

\textbf{Mechanism}: the same paradigm structure that produces the §6.2.2
cross-engine divergence at saturation also explains the within-engine
convergence. Once origin-edge insertion-refusal dominates the dynamics
(the bottleneck behavior controlling network capacity), both SUMO micro
and SUMO meso model that behavior the same way --- both refuse to insert
when the destination edge is at capacity. Lane-level dynamics that micro
adds (sub-link car-following, lane changes, gap acceptance) become
second-order to the insertion-bound queueing. The two paradigms are
interpretively distinct in the free-flow and moderately-congested
regimes; they collapse to the same controlling dynamics under
saturation.

Trip-completion counts confirm the saturation-bound interpretation: SUMO
micro completes 123,735 trips at 200K (61.9 \%) vs SUMO
meso\textquotesingle s 116,270 (58.1 \%) --- micro completes slightly
\emph{more} despite the slower per-trip resolution, because micro models
finer headway packing at the insertion bottleneck.
MATSim\textquotesingle s queue-hold paradigm (no insertion-refusal at
network capacity) completes all 200K (100 \%), the same paradigm
asymmetry that §6.2.2 documents.

\textbf{Significance}: SimForge surfaces \emph{two distinct
paradigm-spread phenomena} at saturation density:

\begin{enumerate}
\def\labelenumi{\arabic{enumi}.}
\tightlist
\item
  \textbf{Cross-engine} (§6.2.2): SUMO insertion-refusal ↔ MATSim
  queue-hold produces a 35.5 \% mean-TT divergence at chicago\_200k,
  96.3 \% at nyc\_500k.
\item
  \textbf{Within-engine} (§6.2.4, this section): SUMO micro ↔ SUMO meso
  mean-TT gap shrinks from +27 \% at 1K to +7.8 \% at 200K as both
  resolutions become bound by the same insertion-refusal dynamics.
\end{enumerate}

Both findings emerge from the same underlying mechanism
(insertion-refusal vs queue-hold at network capacity), but they manifest
in different cross-cuts of the engine × mode matrix. The cross-engine
finding is paradigm-vs-paradigm; the within-engine finding is
resolution-vs-resolution within a single paradigm. Their sibling
structure validates the fairness contract\textquotesingle s role
(§6.2.5): both are \emph{interpretable} as paradigm effects only because
Q1-Q3 PASSes on every cell, ruling out input asymmetry.

\textbf{Practical implication for cross-simulator benchmarking}: at
saturation, \textbf{paradigm choice (insertion-refusal vs queue-hold)
dominates within-engine resolution choice}. The 124 × wall premium for
SUMO micro at chicago\_200k purchases a 7.8 \% mean-TT delta --- a poor
fidelity-cost trade at this regime unless the use case specifically
requires lane-level dynamics that mean-TT comparisons do not surface.
The flat-fidelity / steep-wall trajectory is the within-engine analog of
§6.2.2\textquotesingle s flat-fidelity-then-paradigm-snap trajectory in
the cross-engine direction.

\subsection{The fairness contract as the connecting
tissue}\label{625-the-fairness-contract-as-the-connecting-tissue}

All four of the above findings share a structural property: each is
\emph{interpretable} only because the fairness contract (Q1
byte-identity, Q2 same SCC, Q3 same trip count target) is empirically
demonstrated. Without Q1-Q3, the BFS speedup numbers could be attributed
to a silently-changing input set; the paradigm-divergence and
within-engine convergence findings could be attributed to one engine
receiving easier OD pairs; the cross-platform finding could be
attributed to subtle version drift in the canonical bundle. With Q1-Q3
PASS at every tier on every cell, each of the four findings is
\emph{necessarily} attributable to the mechanism named.

The fairness contract is therefore not just a C1 implementation detail
but the \emph{methodological substrate} that lets the four empirical
findings stand. This generalizes: any cross-simulator benchmarking
framework that does not establish input-byte-identity before reporting
cross-engine numbers is leaving an attribution gap that subsequent
findings cannot cleanly close.

\begin{center}\rule{0.5\linewidth}{0.5pt}\end{center}

\section{Relationship to prior
work}\label{63-relationship-to-prior-work}

\subsection{Cross-simulator benchmarking
efforts}\label{631-cross-simulator-benchmarking-efforts}

Prior cross-simulator studies in the urban-mobility-simulation
literature fall into three rough categories: (a) single-simulator
calibration studies that report fidelity vs observed data within one
engine, (b) two-simulator bake-offs that compare a small number of
metrics on a single shared scenario, and (c) tool-comparison surveys
that catalog feature differences without per-scenario measurement.
SimForge contributes to a category prior work has only thinly populated:
(d) a \emph{reproducible} multi-simulator benchmark framework with
hash-pinned inputs, byte-deterministic adapters, and a formal
cross-engine fairness contract.

The closest analogs in adjacent computational fields are the MLPerf
benchmark suite for machine learning (cited in plan §2.6), the SPEC
benchmark family for system performance, and ReproZip for
computational-experiment packaging. SimForge\textquotesingle s design
draws explicit inspiration from MLPerf\textquotesingle s structured
reproducibility checklist and from the FAIR (Findable, Accessible,
Interoperable, Reusable) data principles, both cited in §2.6 of the
original plan. The SimForge canonical schema + adapter pattern +
fairness audit are the transportation-simulation analogs of
MLPerf\textquotesingle s reference implementations + closed-division
submission rules + accuracy gates.

\subsection{Reproducibility in HPC and scientific
computing}\label{632-reproducibility-in-hpc-and-scientific-computing}

The cross-platform reproducibility finding (§5.6.3, §6.2.3) connects to
a well-documented but rarely-measured phenomenon in numerical computing:
identical version numbers do not guarantee bit-identical outputs across
platforms. The plan cited this in §2.7 ("non- associative
floating-point arithmetic, nondeterministic thread scheduling, divergent
random-number streams, and compiler-dependent optimizations"). The
empirical SimForge measurement contributes a specific data point to that
literature: a 2.95 \% MATSim mean-TT shift between two JVM builds of the
same Java major version, with no other variable changed.

This data point is small in any single dimension but cumulatively
relevant: cross-simulator benchmarking frameworks that report numbers
without specifying the execution context\textquotesingle s binary
closure are reporting results that are not bit-reproducible by
definition. The SimForge container (Wave 2) closes this gap
operationally; the §5.6.3 measurement quantifies what was previously
hand-waved.

\subsection{The engine roster
decision}\label{633-the-engine-roster-decision}

The plan §1.10 named five engines (SUMO, MATSim, POLARIS, LPSim,
QarSUMO). The shipped framework integrates three (SUMO, MATSim,
DTALite). The substitution decision is documented at
\texttt{doc/engines/THIRD\_ENGINE\_OPTIONS.md}. The relevant prior-work
context: at the time of the substitution decision (Phase B, 2026-04-27),
the most-recent published cross-simulator studies in the transportation
literature also typically integrated 2-3 engines, not 5. The shipped
roster is therefore in the median of comparable work, while the
documented rule-out chain (LPSim integration attempted and abandoned
with full retrospective; QarSUMO source unavailable; POLARIS
license-gated; CityFlow effectively abandoned) contributes to the
literature on \emph{what engines are practically integrable} in a
2026-era research codebase.

\begin{center}\rule{0.5\linewidth}{0.5pt}\end{center}

\section{Limitations and threats to
validity}\label{64-limitations-and-threats-to-validity}

The threats catalogued in §4.6 (random seed, JVM warm-up, trip-count
asymmetry, OS scheduling noise) manifested as documented in §5.7. The
limitations below are \emph{broader} --- they bound what claims SimForge
can make even with all the within-scope work delivered.

\subsection{Single-context fidelity vs reality is not
measured}\label{641-single-context-fidelity-vs-reality-is-not-measured}

SimForge measures \emph{inter-simulator agreement} (Q1-Q4 audit, Tables
5.1 and 5.2), not \emph{fidelity vs observed traffic}. The
\texttt{evaluation/metrics/fidelity.py} module implements RMSE, GEH, and
KS statistics, but the ingestion pipeline for Chicago Traffic Counts,
LADOT, or NYC DOT loop-detector data is not built. The framework can in
principle support fidelity-vs-reality measurement; the data acquisition
+ sensor-to-network spatial join + temporal alignment work was out of
scope.

The shipped substitute --- cross-engine TT comparison framed as paradigm
divergence (§5.6.2) --- is informative about engines but silent about
reality. A future researcher who wired in observed baselines would
extend SimForge\textquotesingle s reach from "are the engines
self-consistent?" to "which engine is closer to reality on this
scenario?"

\subsection{The 5M-trip tier was not
run}\label{642-the-5m-trip-tier-was-not-run}

Plan §1.13 declared three demand tiers (50 K, 500 K, 5 M trips per
city). The shipped tier ladder covers 1 K, 10 K, 50 K, 200 K, and 500 K
(max: nyc\_500k\_car). The 5 M tier was not run; even with Phase
14\textquotesingle s BFS speedup, the projected cold-cache wall on
Cardinal for chicago\_5M\_car exceeds practical thesis-window compute
budgets. Pre-Phase-14 the gap was structural (\textasciitilde600 h
projected); post-Phase-14 it is budgetary (\textasciitilde50-100 h
projected, requiring multi-day uninterrupted SLURM allocations). The 500
K tier provides a 5 × scale step over the plan\textquotesingle s
middle tier (50 K → 500 K) and is sufficient for the scalability
analysis the plan scoped.

\subsection{DTALite is missing at the 50 K +
tier}\label{643-dtalite-is-missing-at-the-50-k--tier}

The path4gmns 0.10.0 bundled DTALite C++ binary caps internal OpenMP at
4 threads regardless of allocation (Phase 12.5 discovery, 2026-05-03).
The la\_50k\_car DTALite cells timeout at the 4 h budget; extrapolation
to 200 K + would require \textasciitilde100-250 h per seed.
DTALite\textquotesingle s contribution to the cross-engine matrix is
therefore confined to the 1 K and 10 K tiers, where it adds the UE/DTA
paradigm to the comparison. At 50 K + the matrix reports SUMO + MATSim
only.

Mitigation paths exist (replace bundled binary with a multi-thread
build; evaluate \texttt{path4gmns.find\_ue} pure-Python solver;
integrate alternative DTA tool) and are catalogued in §6.5 below. None
are in scope for the shipped framework.

\subsection{N = 5 seeds, not N =
10}\label{644-n--5-seeds-not-n--10}

Plan §3.4 declared \emph{"an additional N = 10 repeated runs"} per
tuple. The shipped matrix uses N = 5. At N = 5 the
Student\textquotesingle s-t 95 \% confidence intervals are wider than
they would be at N = 10, but the R = 1 − σ/μ measurements are already at
0.95-1.00 across all non-degenerate cells, indicating the variance is
sufficiently small that doubling N would not change any reported claim.
The trade preserves the reproducibility intent at half the compute
budget.

\subsection{No grid-search
calibration}\label{645-no-grid-search-calibration}

Plan §3.4 declared a calibration protocol --- fixed grid search of
3-5 settings per parameter per city, selection by lowest RMSE subject to
throughput ≥ real-time/4. SimForge ships \emph{no calibration}: each
adapter uses default engine config sourced from
\texttt{adapters/common/vehicle\_types.py} and per-engine config
templates. The rationale is fairness --- per-engine tuning would
optimize one engine for an observable (e.g., link counts) that the
cross-engine comparison is supposed to measure neutrally, biasing the
framework\textquotesingle s core claim. The fairness contract trades
calibration latitude for attribution clarity.

\subsection{Code-version drift between time-separated
measurements}\label{646-code-version-drift-between-time-separated-measurements}

The §5.6.3 original measurement of a 2.95 \% MATSim shift was
reinterpreted (§5.6.3.3) as code-version drift between Pitzer Phase 12
(2026-05-02) and Cardinal Phase 14.13 (2026-05-20), not as
cross-platform variability. At a fixed code version, Chapter 5 Tables
5.1 + 5.2 + §5.6.2 numbers are reproducible bit-identically on any
x86\_64 system (§5.6.3.1) and to 4-sig-fig agreement on Mac arm64
(§5.6.3.2). The remaining limitation is that any future benchmark adding
new commits to the routing layer (canonical\_routes BFS, adapter glue,
bundle generator) may produce up to \textasciitilde3 \% mean-TT drift
even on the same machine + bundle. The pinned-digest container is the
operational mechanism that prevents this drift; a future replicator who
pulls \texttt{ghcr.io/phanidharakula/simforge:db8d786} gets the exact
Chapter 5 code, not whatever the active dev branch happens to be.

\subsection{Container mode is opt-in, not
default}\label{647-container-mode-is-opt-in-not-default}

The canonical SBATCH wrappers default to host venv execution; users must
explicitly set \texttt{SIMFORGE\_USE\_CONTAINER=1} to switch into the
pinned-digest container. The decision preserves backward compatibility
for the existing benchmark workflow, but means the shipped Chapter 5
numbers were measured in the host venv context. A future thesis re-pin
could flip the default once the container is the canonical reference (a
possible deliverable for a v1.1 release).

\begin{center}\rule{0.5\linewidth}{0.5pt}\end{center}

\section{Future work}\label{65-future-work}

The following directions are concrete and bounded --- each is described
with the specific blocker that prevented inclusion in the shipped
framework and the path to closure.

\subsection{Observed-baseline
fidelity}\label{651-observed-baseline-fidelity}

\textbf{Goal}: wire RMSE / GEH / KS metrics
(\texttt{evaluation/metrics/fidelity.py}) to observed link counts and
trip-time distributions from CDOT, LADOT, and NYC DOT loop-detector
data. Output: per-corridor + network-level fidelity tables for each
engine, framed as "how close to reality, by city × engine".

\textbf{Effort}: \textasciitilde2-3 weeks. Includes data acquisition
(CDOT/LADOT/NYC DOT loop-detector exports), sensor-to-canonical-network
spatial join (sensor lat/lon → OSM-way → canonical link ID), temporal
alignment to the scenario\textquotesingle s simulated hour, RMSE/GEH/KS
computation per corridor.

\textbf{Value}: extends SimForge\textquotesingle s claim from "engines
are self-consistent" to "engines are X-percent away from reality on this
scenario". Critical for any practitioner using simulation for planning
decisions.

\subsection{DTALite native rebuild for the 50 K +
tier}\label{652-dtalite-native-rebuild-for-the-50-k--tier}

\textbf{Goal}: replace the path4gmns 0.10.0 bundled DTALite C++ binary
(4-thread OpenMP cap) with a manually-compiled DTALite that respects
\texttt{OMP\_NUM\_THREADS}. Restores DTALite participation in the 50 K,
200 K, and 500 K tier benchmarks.

\textbf{Effort}: \textasciitilde1 week. Includes upstream DTALite C++
source build on the target platform (CMake, OpenMP, GMNS data format
support), SimForge adapter modification to call the rebuilt binary
instead of the bundled one, runtime verification on la\_50k\_car
(current timeout target).

\textbf{Alternative}: evaluate \texttt{path4gmns.find\_ue} (pure-Python
UE solver). Slower per iteration but unconstrained by the bundled
binary\textquotesingle s thread cap. Either path closes the same gap.

\textbf{Value}: extends the cross-engine matrix at 50 K + to include all
three paradigms (microscopic queue, activity-based queue, DTA
equilibrium) rather than two.

\subsection{GTFS transit demand
integration}\label{653-gtfs-transit-demand-integration}

\textbf{Goal}: extend the canonical demand schema from car-only to
multimodal by integrating GTFS (General Transit Feed Specification)
schedules for the bus + rail subset of trips. The schema already
supports \texttt{mode\ ∈\ \{car,\ transit,\ bike,\ walk\}}; the missing
piece is the SUMO PT + MATSim transit-routing adapter wiring.

\textbf{Effort}: \textasciitilde3-4 weeks. SUMO PT requires
\texttt{\textless{}busStop\textgreater{}} /
\texttt{\textless{}trainStop\textgreater{}} additionals,
\texttt{ptlines.xml} schedules, and
\texttt{\textless{}vType\textgreater{}} with \texttt{vClass="bus"}.
MATSim transit requires \texttt{transitSchedule.xml} and the qsim
\texttt{mainMode=car,bus,rail} config. Both adapters need PT-specific
output parsers.

\textbf{Value}: closes plan §1.10 Objective 2 "GTFS" alignment that
was dropped (D8 in PLAN\_DEVIATIONS.md), opens cross-engine
comparison on the transit-mode subset.

\subsection{POLARIS integration if license becomes
accessible}\label{654-polaris-integration-if-license-becomes-accessible}

\textbf{Goal}: integrate POLARIS as a fourth shipping engine if
Argonne\textquotesingle s licensing process becomes navigable for
non-Argonne researchers.

\textbf{Effort}: unknown; the integration itself is \textasciitilde3
weeks of adapter work plus the license-acquisition cycle (plan §1.11
documents the license as "freely available for academic research via
request form", but the request-to-grant cycle is opaque to outsiders).

\textbf{Value}: extends the cross-engine matrix to four paradigms,
strengthens the plan-§1.11 C2 commitment by partial recovery of the
original 5-engine roster.

\subsection{5 M-trip tier}\label{655-5-m-trip-tier}

\textbf{Goal}: extend the demand tier ladder to chicago\_5M\_car or
nyc\_5M\_car for the scalability analysis.

\textbf{Effort}: bounded by compute budget. With Phase 14 cold-cache
wall at 0.5-0.7 trips/s/worker (empirical from §5.6.3), the BFS pass for
5 M trips on a metropolitan-scale SCC projects to 100-200 h on a
24-worker Cardinal node. The engine sims after BFS would land in 5-10 h.
Total \textasciitilde110-210 h per cold-cache run, requiring multi-day
uninterrupted SLURM allocations.

\textbf{Value}: completes the original plan §1.13 scope. The 500 K
tier (max shipped) already provides a 50 × scale step over 10 K, so the
additional 10 × from 500 K → 5 M is incremental confirmation, not new
qualitative information.

\subsection{Cross-machine container
verification}\label{656-cross-machine-container-verification}

\textbf{Goal}: verify the container produces bit-identical outputs
across multiple machines (Cardinal, AWS, Azure, a
colleague\textquotesingle s cluster), not just within Cardinal.

\textbf{Effort}: \textasciitilde1 week. Requires access to ≥ 2
non-Cardinal compute environments, pull the same container digest, run
the same benchmark on each, diff the outputs. The expectation (per
§5.6.3) is that container-mode results across machines are
bit-identical, since the container\textquotesingle s binary closure is
platform-independent.

\textbf{Value}: empirically substantiates the §5.6.3 claim that the
container is \emph{the} canonical reproducibility target across
machines.

\subsection{Zenodo DOI minting and JOSS
publication}\label{657-zenodo-doi-minting-and-joss-publication}

\textbf{Goal}: publish the SimForge artefact as a citable software paper
(Journal of Open Source Software) with a Zenodo DOI for stable citation.
Closes the plan §1.10 Objective 4 "Open Test Suite Release"
commitment with a peer-reviewed citation.

\textbf{Effort}: \textasciitilde1 week for JOSS paper draft + Zenodo DOI
minting + public-repo flip + advisor sign-off.

\textbf{Value}: stable citation handle for SimForge for any future paper
or thesis that builds on the framework.

\begin{center}\rule{0.5\linewidth}{0.5pt}\end{center}

\section{Conclusion}\label{66-conclusion}

This thesis presented SimForge, a reproducible cross-simulator testing
framework for urban mobility simulation. The four plan contributions
(canonical schema, deterministic adapters, pinned-digest execution,
hardware-normalized KPIs) are delivered, with documented adjustments to
the engine roster and the per-watt + calibration mechanisms. The
framework integrates SUMO, MATSim, and DTALite, covering three distinct
simulation paradigms --- microscopic queue, activity-based queue, and
user-equilibrium DTA --- under a single canonical schema with
byte-deterministic adapters and a programmatic cross-engine fairness
audit.

Three empirical findings emerged from the build that contribute to the
cross-simulator benchmarking literature beyond the original C-row
deliverables:

\begin{enumerate}
\def\labelenumi{\arabic{enumi}.}
\item
  \textbf{Phase 14 BFS deduplication} reduced the chicago\_200k\_car
  benchmark wall by approximately 20 × cold-vs-cold (141.87 h → 7.14 h)
  and 228 × for warm-cache re-runs. The 500 K-trip tier (nyc\_500k\_car)
  became tractable for the first time, completing in 12 h 8 min wall vs
  a pre-Phase-14 projection of \textasciitilde600 h.
\item
  \textbf{Q4 paradigm divergence at scale} showed that the SUMO meso vs
  MATSim qsim mean-TT ratio is regime-dependent, not regime-stable:
  small-tier convergence (within ±5 \% at 50 K trips) reverses sharply
  at saturation (96.3 \% gap at 500 K trips on NYC). The
  convergence-then-divergence pattern is attributable to paradigm-level
  mobsim choices (SUMO\textquotesingle s insertion-refusal vs
  MATSim\textquotesingle s queue-hold), and is interpretable as such
  only because the fairness contract (Q1-Q3 PASS) rules out input
  asymmetry as a cause.
\item
  \textbf{Cross-platform reproducibility limits} were measured at
  approximately 0.2-3 \% per engine, with the largest shift being a 2.95
  \% MATSim mean-TT divergence between brew openjdk@17 and Debian
  openjdk-17-jre-headless, even with the same MATSim JAR and the same
  \texttt{lastIteration=0} / \texttt{numberOfThreads=1} configuration.
  The pinned-digest container delivered in Wave 2 is the operational
  mechanism that closes this gap, providing bit-identical
  reproducibility across machines as the canonical citation target.
\end{enumerate}

The fairness contract (Q1 byte-identity of feasibility verdicts, Q2 same
SCC network, Q3 same trip-count target) is the methodological substrate
that makes each of the four findings interpretable. Without it, the
Phase 14 speedup measurements could be attributed to silently changing
input; the cross-engine and within-engine paradigm findings could be
attributed to engines (or modes) receiving different OD pairs; the
cross-platform shift could be attributed to canonical-bundle drift. With
Q1-Q3 demonstrated PASS at every tier on every shipped cell, each
finding is necessarily attributable to its named mechanism.

SimForge is now a defensible \emph{reproducible execution target} for
cross-simulator urban-mobility benchmarks, with the artefact
distribution mechanism (container at
\texttt{ghcr.io/phanidharakula/simforge:db8d786}), the canonical bundle
schema (\texttt{scenarios/\textless{}scenario\_id\textgreater{}/}), and
the audit + scorecard tooling (\texttt{evaluation/audit\_fairness.py},
\texttt{tools/generate\_scorecard.py}) in place. Future researchers
building on this work can pull the container, bind-mount the published
bundles

\begin{itemize}
\tightlist
\item
  hash-pinned OSM data, and obtain bit-identical results to the
  measurements reported in Chapter 5 --- modulo the documented ±3 \%
  cross-context shift that itself is part of the
  framework\textquotesingle s empirical contribution.
\end{itemize}

The framework\textquotesingle s value is not the specific numbers it
produced for chicago\_200k\_car or nyc\_500k\_car, but the
\emph{mechanism} by which any future cross-simulator comparison can be
made fair, deterministic, and bit-reproducible. The
plan\textquotesingle s central thesis --- that \emph{cross-simulator
benchmarking requires a reproducible, hardware-normalized,
fairness-audited framework, and that such a framework can be built} ---
is supported by both the artefact and the empirical findings that
emerged from building it.

\begin{center}\rule{0.5\linewidth}{0.5pt}\end{center}

\section{Reproducing this chapter\textquotesingle s
claims}\label{67-reproducing-this-chapters-claims}

Every empirical claim in this chapter traces to a specific artefact that
a future reader can re-verify:

\begin{longtable}[]{@{}lll@{}}
\toprule\noalign{}
Claim & Artefact & Verification command \\
\midrule\noalign{}
\endhead
\bottomrule\noalign{}
\endlastfoot
C1 schema validates & \texttt{pipeline/validation/validate\_bundle.py} &
\texttt{python\ -m\ pipeline.validation.validate\_bundle\ scenarios/chicago\_1k\_car} \\
C2 byte-deterministic adapters &
\texttt{tests/test\_adapter\_determinism.py} &
\texttt{python\ -m\ pytest\ -m\ determinism} \\
C3 container reproducibility &
\texttt{ghcr.io/phanidharakula/simforge:db8d786} &
\texttt{singularity\ pull\ docker://ghcr.io/phanidharakula/simforge:db8d786\ \&\&\ apptainer\ exec\ -\/-pwd\ /workspace/SimForge\ containers/simforge\_db8d786.sif\ python\ -c\ "import\ adapters.sumo.sumo\_adapter;\ print(\textquotesingle{}OK\textquotesingle{})"} \\
C4 fairness audit + scorecard & \texttt{evaluation/audit\_fairness.py},
\texttt{tools/generate\_scorecard.py} &
\texttt{python\ -m\ evaluation.audit\_fairness\ runs/benchmark\_small\ \&\&\ python\ -m\ tools.generate\_scorecard\ runs/benchmark\_small} \\
Phase 14 \textasciitilde20 × cold-vs-cold speedup &
\texttt{doc/EXPERIMENT\_LOG.md} 2026-05-18 + 2026-05-19 entries &
Compare Cardinal job 9332478 (Phase 13) vs 9971041 (Phase 14) wall
times \\
Phase 14 nyc\_500k 12 h cold-cache run & Cardinal job 9971042 &
\texttt{runs/benchmark\_large/nyc\_500k\_car/benchmark\_results\_benchmark\_large.json}
\texttt{started\_at} / \texttt{completed\_at} \\
Q4 paradigm divergence at nyc\_500k\_car & §5.6.2 + audit\_fairness.txt
&
\texttt{python\ -m\ evaluation.audit\_fairness\ runs/benchmark\_large/nyc\_500k\_car} \\
Cross-platform 2.95 \% MATSim shift & §5.6.3 + 2026-05-20
EXPERIMENT\_LOG entry & Diff
\texttt{runs/container\_smoke/chicago\_1k\_car/} vs
\texttt{runs/benchmark\_small/chicago\_1k\_car/} \\
Fairness contract Q1-Q3 PASS at every tier &
\texttt{audit\_fairness.txt} per run & Run the audit command above on
each run directory \\
\end{longtable}

The full reproduction recipe for the entire thesis is at
\href{../REPRODUCING.md}{\texttt{doc/REPRODUCING.md}}. The
container-based reproduction path (the recommended one) is at
\href{../CONTAINER_USAGE.md}{\texttt{doc/CONTAINER\_USAGE.md}}.
